
%%%%%%%%%%%%%%%%%%%%%%%%%%%%%%%%%%%%%%%%
%           main body
%%%%%%%%%%%%%%%%%%%%%%%%%%%%%%%%%%%%%%%%

%-----------------------------------
%	TITLE PAGE
%-----------------------------------

\title[第十章\\
函数项级数]{\texorpdfstring{\LARGE \bfseries 第十章\quad 函数项级数}{第十章 函数项级数}} % The short title appears at the bottom of every slide, the full title is only on the title page
\author[\smaller \S 10.5\\
多项式逼近] % Your institution as it will appear on the bottom of every slide, maybe shorthand to save space
{\Large \bfseries \S 10.5 用多项式逼近连续函数}
% \author{Singular Value Decomposition} % Your name
% \date{\today} % Date, can be changed to a custom date
\date{}

%------------------------------------------------

\begin{document}

%------------------------------------------------

\begin{frame}

\titlepage % Print the title page as the first slide

\end{frame}

%-----------------------------------
%	PRESENTATION SLIDES
%-----------------------------------

%------------------------------------------------

\begin{frame}
\frametitle{引论}

目的: 用尽可能“简单的”函数来\alert{近似}“复杂函数”

\pause
\vspace{1em}

\begin{block}{问题}
\pause
\begin{itemize}
\item 有哪些函数足够简单? \pause ~ $\rightsquigarrow$ ~ 多项式, 三角函数, $\dots$
\pause
\item \alert{近似}如何定义? \pause ~ $\rightsquigarrow$ ~ 需要定义某种“距离”
\item 是不是任何“复杂函数”都能被“简单函数”近似?
\end{itemize}
\end{block}

\end{frame}

%------------------------------------------------

\begin{frame}
\frametitle{多项式一致逼近}

\begin{Def}[多项式一致逼近]

\end{Def}

\end{frame}

%------------------------------------------------

\begin{frame}
\frametitle{Weierstraß 第一逼近定理}

\begin{thm}[Weierstraß 第一逼近定理]
设 $f(x) \in C[a, b],$ 则 $\forall~\varepsilon > 0,$ 存在多项式 $P(x)$ 使得 $\forall~x\in [a, b],$ 有 $|P(x) - f(x)| < \varepsilon.$
\end{thm}

\pause

等价地来说, 就是有
$$
d(P, f) := \sup_{x\in[a,b]}|P(x) - f(x)| < \varepsilon.
$$

\pause

Weierstraß 第一逼近定理的重要性在于, 它是 $L^p$ 空间可分性证明的重要组成部分.

\end{frame}

%------------------------------------------------

\begin{frame}
\frametitle{Korovkin 定理}

\end{frame}

%------------------------------------------------

\begin{frame}
\frametitle{Bernstein 定理的证明}

\end{frame}

%------------------------------------------------

\begin{frame}
\frametitle{Weierstraß 第一逼近定理的证明}

\end{frame}

%------------------------------------------------

\begin{frame}
\frametitle{Korovkin 定理的证明}

\end{frame}

%------------------------------------------------

\end{document}
